
\section{The Samaritan}

Luke 10:23–37

“Folkebladet,” September 19, 1917


The last part of today’s Gospel is very well known, namely, the account of the Good Samaritan. Jesus told this event to the man who asked: “What shall I do to inherit eternal life?”

There is, as everyone will see, a considerable amount of Law in Jesus’ answer to this man. And many do not like the preaching of Law, or, if you will, of morality, works. Look and see, and you will be astonished at how often he speaks of works. We are so accustomed to the question of how one is to be saved being answered, “Believe in Jesus,” that we are almost startled when some other answer is given.

“Keep the commandments, and you shall live,” says Jesus. And concerning the day of judgment, of which Jesus speaks in Matthew 25, he says that what is decisive as to who shall be placed at the King’s right hand or at his left is their works, what they have done or neglected to do.

How are we to understand this? Does Jesus teach another way of salvation than, for example, Paul, and than the one ordinarily preached in our services?

Oh no, surely he does not. He himself has said as plainly as possible: “The one who believes in me has eternal life.” “Whoever believes and is baptized shall be saved,” and there are many other such words of his.

The different answers to the question of salvation do not stand in conflict with one another. There is complete agreement between them. For why are we to believe in Jesus? Many will answer this question: We are to believe in Jesus in order to escape perdition—to avoid going to hell. But this answer is hardly the right one. When the angel told Mary that she was to become the mother of the Savior of the world, he said: “You shall call his name Jesus, for he shall save his people from their sins!”

We are therefore to receive Jesus, to believe in him, in order to be set free from sin and enabled to keep God’s commandments.

Paul expresses it thus: God sent his Son and condemned sin in the flesh, *in order that the requirement of the Law might be fulfilled in us.*

There are many who seem to understand matters as though, provided only that one believes in Jesus and comforts oneself with his atoning death, one will be admitted into the blessedness of heaven, whether one has been willing to keep the commandments, to love God and to love one’s neighbor, or not. But this is self-deception, and it is against this kind of self-deception that Jesus so often and so insistently speaks.

In heaven or on earth there is no blessedness to be attained without goodness: “Blessed are the pure in heart, for they shall see God,” says Jesus. And Jesus will give us that pure heart when we give ourselves over to him without reserve, body and soul. And this is to believe in him:

“The faith that Jesus embraces
and that can cleanse the heart,
the faith by which one reaches
heaven’s land of joy,
that faith is strong, and must
bring all the world beneath it,
in sorrow and in gladness
be able to endure its trial!”

These two who passed by the fallen and wounded man, although they *saw him* and the distress he was in, were doubtless both well satisfied in the consciousness that they possessed “the right faith.”

What, then, was it that they lacked, and that Jesus says we must have if it is to be well with us eternally?

It was the disposition that has compassion upon distress and desires to relieve it, the disposition that has compassion *because it loves.* The Good Samaritan possessed that disposition, and therefore he helped the man who had fallen among robbers. The Samaritan’s love, goodness of heart, and compassion drove him to show mercy toward the fallen man although he was a Jew, a political and religious opponent.

As Jesus said to the man who asked for the way to life, “Go, and do likewise,” so he says to us. We never attain the goal of our life, which is *God himself,* by having the right understanding of Jesus’ saving work and holding that understanding to be true. Many think so, but are grievously mistaken. It does not avail for salvation unless it is able to *change our disposition and our life.*

The goal of our life is attained only by our being transformed more and more to *become like God,* becoming merciful through that love and power in the heart which “is poured out into our hearts through the Holy Spirit who has been given to us.” The question that, according to today’s Gospel, is so important is not this: “Can I have a blessed death?” But this: “*Do I have the Samaritan disposition?*”

There are sincere souls who become afraid of themselves at this question and who, with shame and sorrow, must answer: I have so little of this disposition. And such persons are not *satisfied* to have so little. They long for more.

How can we obtain more of this disposition? I will answer: *By becoming glad.* Become more glad, and you will become more merciful, more good. And how can I become more glad?

There is, indeed, so much here that is fitted to make us sorrowful: anxiety for ourselves, anxiety for those dear to us, anxiety for the future, and more of the same kind.

Oh, my brother and sister, *God* can make you glad. He does it *through his Word:*

“Whatever here in the world may grieve my soul,
with *Jesus and the Word* it always turns out well.
I do not always feel and know my happiness;
yet in faith upon the Word my soul has peace.
Oh, blessed rest:
to believe in Jesus
and have God’s Word dwelling within!”

“When life grows dark, when anxious I and weak,
I hasten to the Word, and there it becomes day.
When others are rich and I am not as they,
I sit within the Word, and *there I am rich.*
Oh, blessed rest:
to believe in Jesus
and have God’s Word dwelling within!”

In his Word the Lord says to us: Go to Christ as you are—as helpless and sinful as you feel yourself to be—and you receive righteousness and peace with God, and *then you become glad—blessed.* “Blessed is the one whose transgression is forgiven, whose sin is covered.”

But this deep peace of God in Christ, which surpasses all understanding, is not a peace in which to idle, a peace of *inactivity.* No, it is working peace, peace for working with willingness and zeal, working in our heavenly Father’s great vineyard. We do not measure our children’s love for us first and chiefly by their words about how dearly they love father and mother, but by *their willingness to serve and obey us.* So does God also.

“We have been created a second time, in Christ, for good works which God himself has prepared for us, that we should walk in them.” And there is enough to do in our heavenly Father’s vineyard. Enough for everyone. There is the foreign mission. And there is the home mission. So many of our brothers and sisters in the world have “fallen among robbers.” And many of them you and I see. Shall we do as the two did: *pass by?* Or shall we *have compassion and help,* even though this costs both *trouble* and—*money?*

Oh, my brother and sister, you who have had your feet set upon the rock and who are permitted to go about here in this world of sorrow with the blessed assurance in your heart of your right and standing as a child of God—*work for the Lord!* You have the opportunity. Use it. Let it be evident in your conduct that you not only look forward to “a blessed death,” but *rejoice in the lifelong work for your own salvation and that of your brothers!*

Do not let your Christianity consist in *having things well,* you and your own, while you close your heart against the great, crying need around you! No, let it *drive you and constrain you to action and sacrifice!* Learn from him who, in love for you, sacrificed *all*—sacrificed every drop of blood, for you, *for me!* Amen.

E. B.
